% !TeX root = ../../main.tex
\subsection{割合の問題}

    \begin{techniquebox}[tech:04]{\textbf{割合を扱うときの鉄則！}}
        \textbf{割合は\fitblankbf[分数]{}や\fitblankbf[小数]{}に直して
        \fitblankbf[かけ算]{}で扱う！}
    \end{techniquebox}

    割合の問題のうち、代表的なものに、\uwave{「生徒数の問題」}と
    \uwave{「食塩水の問題」}がある。それぞれについて、深掘りしていこう。

    \subsubsection{生徒数の問題}
    \vspace{1em}

        \begin{techniquebox}[tech:05]{\textbf{生徒数の問題}}
            \textbf{(男子の生徒数) + (女子の生徒数) = (全校生徒数)} という考え方を使う!
        \end{techniquebox}

        Technique \ref{tech:05} は最低限使えて欲しい考え方ではあるが、余裕があれば、
        \uwave{生徒数の増減に着目した}考え方も使えて欲しい。

        \[\textbf{(男子の増減) + (女子の増減) = (その年の生徒数の増減)}\]

        という式を立てることができると計算が楽になる。

        \begin{example}[【例題 3】]
            （ここに過不足の問題の例題文を書く）
        \end{example}

    \subsubsection{食塩水の問題}
    \vspace{1em}

        \begin{techniquebox}[tech:06]{\textbf{食塩水の問題}}
            \textbf{(混ぜる前の塩の量) = (混ぜた後の塩の量)} という考え方を使う!
        \end{techniquebox}

        塩の量は、(食塩水の量) $\times$ (塩の濃度) で求められるので、Technique \ref{tech:04} より、

        \[\textbf{(塩の量) = (食塩水の量)} \times \left(\dfrac{\text{\textbf{濃度}}}{\bm{100}}\right)\]

        \begin{example}[【例題 4】]
            （ここに過不足の問題の例題文を書く）
        \end{example}

\newpage
